% ============================================================
%  Деталі реалізації — вставити ПЕРЕД секцією
%  "Реалізація та експериментальна верифікація"
% ============================================================

\subsection{Інструмент запиту до пам'яті (Memory Query Tool)}
\label{sec:memory-query-tool}

Інтерфейсом витягування контексту для агента слугує єдиний MCP"=інструмент \texttt{workflow\_memory\_query}, що забезпечує уніфікований доступ до всіх трьох рівнів пам'яті.
Вибір єдиного інструменту замість окремих точок доступу до кожного рівня є принциповим проєктним рішенням: агент формулює один запит природною мовою, а підсистема пам'яті самостійно маршрутизує його до релевантних рівнів та агрегує відповіді.

Параметри інструменту:
\begin{itemize}[leftmargin=*, nosep]
  \item \texttt{task\_description} (рядок, обов'язковий) --- природномовний опис поточної задачі або оперативної ситуації.
  \item \texttt{scope} (перелічення: \texttt{all}\,$|$\,\texttt{domain}\,$|$\,\texttt{workflow}\,$|$\,\texttt{practitioner}) --- рівні пам'яті для запиту. За замовчуванням: \texttt{all}.
  \item \texttt{token\_budget} (ціле число, за замовчуванням 8\,000) --- максимальна кількість токенів у зібраному дайджесті.
  \item \texttt{filters} (об'єкт, опціональний) --- структуровані фільтри: компонент, рівень критичності, часовий діапазон.
\end{itemize}

Інструмент повертає \emph{структурований прозовий дайджест} із такими складовими:
\begin{enumerate}[nosep]
  \item заголовки секцій відповідно до кожного рівня пам'яті, з якого витягнуто записи;
  \item оцінка релевантності (relevance score) для кожного запису;
  \item посилання на джерело (провенанс) --- конкретний коміт, рішення командира, запис трекера або документ домену.
\end{enumerate}

Формат дайджесту спроєктовано для безпосередньої ін'єкції у контекстне вікно агента без додаткової постобробки.
Токенний бюджет забезпечує контрольовану вартість ініціалізації: підсистема ранжує записи за релевантністю і відсікає нижні записи, коли сумарний обсяг перевищує бюджет.

\textbf{Військова аналогія.}
У системах управління БПЛА або ситуаційних центрах еквівалентом є \emph{уніфікований інтерфейс розвідувального запиту}: оператор формулює інформаційну потребу один раз, а система агрегує дані з усіх доступних джерел --- радарів, БПЛА"=розвідки, агентурних даних, супутникових знімків --- у єдину розвідувальну зведення.
Альтернатива --- ручний послідовний запит до кожного сенсора окремо --- призводить до затримок, фрагментації інформації та підвищеного когнітивного навантаження на оператора.
Уніфікований інтерфейс запиту є реалізацією принципу \emph{інформаційного злиття} (information fusion), що є стандартним елементом архітектур C4ISR~(командування, управління, зв'язок, комп'ютери, розвідка, спостереження, рекогносцировка).

\subsection{Оркестратор довгострокових задач (Long-Term Task Orchestrator)}
\label{sec:orchestrator}

Push"=режим реалізований як окремий Docker"=контейнер з власним cron"=розкладом, ізольований від основного серверу обробки запитів.
Архітектурна ізоляція забезпечує дві властивості: (1)~фонові процеси оновлення не конкурують за ресурси з активними сесіями операторів; (2)~збій оркестратора не впливає на доступність pull"=режиму.

Цикл роботи оркестратора:
\begin{enumerate}[nosep]
  \item \textbf{Виявлення задач.} Оркестратор надсилає запит до трекера задач (Plane API) і отримує список задач із позначкою \texttt{LONG-TERM}.
  \item \textbf{Обчислення пріоритету оновлення.} Для кожної задачі обчислюється пріоритет оновлення за формулою~(\ref{eq:refresh}), яка є монотонно зростаючою функцією трьох сигналів: кількості комітів, що стосуються задачі, кількості коментарів до задачі та часу з моменту останнього pull"=запиту.
  Частота оновлення обмежена знизу одним оновленням на 7 днів та зверху одним оновленням на 24 години для кожної задачі.
  \item \textbf{Генерація виходів.} Для кожної задачі, що підлягає оновленню, оркестратор генерує три структуровані виходи за допомогою LLM"=сумаризатора:
  \begin{itemize}[nosep]
    \item \textbf{Резюме для командира} (executive summary) --- стислий опис змін, що відбулися з моменту останнього оновлення: нові рішення, зміни обстановки, результати виконання суміжних задач.
    \item \textbf{Дельта лінійки інструментів} (tool lineage delta) --- створені, модифіковані та виведені з експлуатації інструменти з відмінностями у схемах параметрів.
    \item \textbf{Відкриті питання} (open questions) --- невирішені проблеми або точки прийняття рішень, що потребують уваги оператора при поверненні до задачі.
  \end{itemize}
  \item \textbf{Індексація та зберігання.} Згенеровані виходи вбудовуються у вектори та зберігаються у колекції Qdrant робочого рівня пам'яті з метаданною позначкою \texttt{source:\,push-refresh}.
  При наступному pull"=запиті до оновленої задачі ці записи автоматично потрапляють у дайджест поряд із записами основної пам'яті, забезпечуючи сесію актуальним контекстом.
\end{enumerate}

\textbf{Військова аналогія.}
Оркестратор є програмним аналогом \emph{автоматизованої підготовки розвідувальних зведень між змінами}.
У штабній практиці черговий офіцер між змінами готує бойове зведення: збирає інформацію з різних джерел, виділяє зміни обстановки, фіксує невирішені питання.
Наступна зміна починає роботу не з порожнього аркуша, а з актуального зведення.
Оркестратор автоматизує саме цей процес: між сесіями операторів він збирає зміни, генерує структуроване зведення і зберігає його у підсистемі пам'яті так, що наступний оператор отримує актуальний контекст автоматично при ініціалізації сесії.
Різниця між ручним та автоматизованим зведенням --- у масштабованості: оркестратор підтримує одночасно десятки довгострокових задач, тоді як людина"=офіцер обмежена когнітивною пропускною здатністю.
